% WM-Concord capstone report
% Adapted from the course template (originally by faculty at Reykjavik University, Iceland)

\documentclass{scrartcl}

% Adapted from an original template by Hlyni Arnórssyni, Reykjavik University, Iceland
%
% ------------------------------ SETTINGS
\usepackage{geometry}

\geometry{
	paper=a4paper, % Paper size
	top=2.5cm, % Top margin
	bottom=2.5cm, % Bottom margin
	left=2.5cm, % Left margin
	right=2.4cm, % Right margin
	headheight=0.75cm, % Header height
	footskip=1.5cm, % Space from the bottom margin to the baseline of the footer
	headsep=0.75cm, % Space from the top margin to the baseline of the header
	%showframe, % Uncomment to show how the type block is set on the page
}

\usepackage{blindtext}
%-------------------------------- Character encoding ----------------------------
\usepackage[T1]{fontenc}
\usepackage[utf8]{inputenc}

%----------------------------- Mathematics packages from AMS ---------------

\usepackage{amsmath, amsfonts, amsthm, amssymb}
\usepackage{braket, nicefrac}

% ----------- International System of Units
\usepackage{siunitx}

%------------------------------ Lists / numbers -------------------------
\usepackage{enumitem, multicol}

%------------------------------- Figure insertions --------------
\usepackage{graphicx, float}  % Use option [H] to force the placement of a figure
\usepackage{keystroke}
\usepackage{pgfplots}\usepgfplotslibrary{units}\pgfplotsset{compat=1.16}




%%%%%%%%%%%%%%%%%%%%%%%%%% Hyperlink References %%%%%%%%%%%%%%%%%%%%%%%%%%%
\usepackage{hyperref}

%--------------------% Storage Path for images %-----------------%
\graphicspath{{graphics/}{Graphics/}{./}}
\usepackage{graphicx,epsfig}
\hypersetup{
   colorlinks   = true,
   urlcolor     = blue,
   linkcolor    = blue,
   citecolor    = red,
   setpagesize  = false,
   linktocpage  = true,
}
\graphicspath{ {fig/} }

\renewenvironment{abstract}{
    \centering
    \textbf{Abstract}
    \vspace{0.5cm}
    \par\itshape
    \begin{minipage}{0.7\linewidth}}{\end{minipage}
    \noindent\ignorespaces
}
% ------------------------------------------------------------------------------------------------------------------------

\begin{document}
\begin{titlepage}
	\centering
	\includegraphics[width=0.6\textwidth]{GW_logo.eps}\par
	\vspace{2cm}
	{\scshape\LARGE Data Science Program \par}
	\vspace{1cm}
	{\scshape\Large Capstone Report - Fall 2026\par}
	\vspace{1.5cm}
	{\huge\bfseries WM-Concord: Do World-Model Diagnostics Agree?\par}
	\vspace{0.5cm}
	{\Large A Cross-Diagnostic Validity Study of State-Equivalence Metrics, Inductive-Bias Probes, and Representation Probing\par}
	\vspace{2cm}
	{\Large\itshape Venkatesh Nagarjuna,\\ Siddardha Reddy,\\ Nikhil Obuleni,\\ Viharika Vemparala\\}\par
	\vspace{1.5cm}
	supervised by\par
	Amir Jafari

	\vfill
	\begin{abstract}
	    The claim that sequence models acquire ``world models'' is tested in the literature with
	    four different diagnostics: next-token validity, Myhill-Nerode state-equivalence metrics,
	    inductive-bias probes, and representation probing. These diagnostics have never been run
	    head-to-head on the same models, and on the same domain they have reached opposite
	    conclusions. We train a controlled zoo of small sequence models on domains with a known
	    ground-truth automaton, run all four diagnostics on every model, and ask whether the
	    diagnostics rank models the same way, why they disagree when they do, and which one best
	    predicts failure on related but shifted tasks. [Results to be added.]
	\end{abstract}
	\vfill
\end{titlepage}
\tableofcontents
\newpage
% ------------------------------------------------------------------------------------------------------------------------
\section{Introduction}

Whether a sequence model trained only on next-token prediction builds an internal model of the
domain that generated its data is a central question for the current foundation-model research
program. The field has produced several diagnostics for answering it, each resting on a different
notion of what ``having a world model'' means. This report describes WM-Concord, a study that
implements four of these diagnostics behind a common interface, applies them to the same set of
models, and measures whether they agree.

% TODO: expand from reports/proposal.md, Section 1.

% ------------------------------------------------------------------------------------------------------------------------
\section{Problem Statement}

Four families of world-model diagnostics exist. Next-token validity checks whether a model only
proposes legal continuations. State-equivalence metrics derived from the Myhill-Nerode theorem
check whether the model treats sequences that reach the same state as equivalent and sequences
that reach different states as distinct. Inductive-bias probes check what a model assumes when
adapted to a new task with very little data. Representation probes check whether the true state
can be decoded from internal activations. These measure different properties, and a model can
pass one while failing another. No published work has run all four on the same models.

We ask four questions. H1: do the four diagnostics induce the same ranking over models? H2:
where they disagree, is the disagreement systematic or noise? H3: which diagnostic best predicts
failure on related-but-shifted tasks? H4: do the findings hold on a held-out domain and a
held-out architecture?

% TODO: expand from reports/proposal.md, Sections 1 and 3.

% ------------------------------------------------------------------------------------------------------------------------
\section{Related Work}

% Draft. Ten or more works required by 2026-10-06. Four are in; six more to add:
% Nanda et al. 2023, Toshniwal et al. 2022 (chess), Hewitt & Liang 2019 (control tasks),
% Belinkov 2022 (probing survey), Hazineh et al. 2023, and one on measurement validity in ML.

\subsection{Representation probing and Othello-GPT}

Li et al.~\cite{li2023othello} trained a GPT-style model on Othello move sequences with no access
to the board or the rules and found that it predicts legal moves with high accuracy. Probes
trained on the residual stream recover the board state, and editing the recovered board
representation changes the model's next-move prediction in the way the edited board implies,
which makes the evidence causal rather than merely correlational. This is the origin of the
representation-probing diagnostic and of the intervention experiment we use to validate it.
Our project differs in that we treat probing as one of four instruments to be compared, not as
the answer, and we pair every probe with a control task so that probe capacity is separated
from genuine encoding.

\subsection{State-equivalence metrics}

Vafa et al.~\cite{vafa2024evaluating} formalize world-model recovery for domains governed by a
deterministic finite automaton and derive two metrics from the Myhill-Nerode theorem: a
compression error, incurred when a model fails to treat two sequences that reach the same state
as equivalent, and a distinction error, incurred when it fails to separate sequences that reach
different states. On taxi routes, Othello, and logic puzzles, models that score well on
next-token accuracy and legal-move rate fail these metrics, and the maps reconstructed from
them contain streets that do not exist. Training on random walks rather than shortest paths
produces more coherent world models. We reimplement these metrics as sampling estimators and
use the authors' detour experiment as the template for our fragility suite. Our contribution
is to measure whether this diagnostic agrees with the other three, which the original paper
does not ask.

\subsection{Inductive-bias probes}

Vafa et al.~\cite{vafa2025inductive} introduce a different instrument: adapt a pretrained model
to small synthetic datasets generated from a postulated world model and check whether its
extrapolations follow that world model. Models trained on orbital trajectories predict
trajectories well but, when adapted to new physics tasks, do not behave as if they know
Newtonian mechanics, and symbolic regression on the adapted models recovers force laws that
are physically meaningless. The conclusion is that sequence-prediction skill does not imply the
world model. We implement this probe as our third diagnostic, including the symbolic-regression
arm with multiple seeds, and compare its verdict against the behavioural and representational
diagnostics on the same models.

\subsection{The disagreement that motivates this study}

Yuan and S{\o}gaard~\cite{yuan2025revisiting} fine-tune seven language models on Othello
sequences, report legal-move accuracy up to 99\% and successful board-state probing on all
seven, and find that the learned representations are similar across architectures, which they
read as evidence for genuine world models. This is the opposite conclusion to
Vafa et al.~\cite{vafa2024evaluating} on the same domain. That two careful studies reach
opposite verdicts with different instruments is exactly the measurement-validity question our
H1 addresses, and the cross-architecture similarity they report bears directly on the
architecture axis of our model zoo.

% ------------------------------------------------------------------------------------------------------------------------
\section{Solution and Methodology}

\subsection{Domains and ground truth}
% TODO: Othello, navigation, interpreter; DFA specification per domain; common generator
% interface emitting (sequence, true state, legal next tokens).

\subsection{Model zoo}
% TODO: transformer, LSTM, Mamba; three scales; three training distributions per domain;
% seeds; run matrix in src/docs/run_matrix.csv; measured throughput.

\subsection{The four diagnostics}
% TODO: D1 next-token validity; D2 compression and distinction estimators; D3 adaptation
% probe and symbolic regression; D4 layer-wise probes with control tasks and activation patching.

\subsection{Fragility suite and pre-registration}
% TODO: per-domain shifted tasks, frozen before diagnostics run.

\subsection{Concordance and predictive-validity analysis}
% TODO: Kendall's tau, Spearman, Kendall's W with bootstrapped CIs; regression of fragility on
% diagnostic scores.

% ------------------------------------------------------------------------------------------------------------------------
\section{Results and Discussion}

\subsection{Reproduction of published results}
% TODO: Othello-GPT probe accuracy and Vafa 2024 navigation metrics, ours beside published.

\subsection{Baselines}
% TODO: legal-move rate per architecture on Othello.

\subsection{Concordance (H1)}
% TODO

\subsection{Sources of disagreement (H2)}
% TODO

\subsection{Predictive validity (H3)}
% TODO

\subsection{Held-out domain and architecture (H4)}
% TODO

% ------------------------------------------------------------------------------------------------------------------------
\section{Discussion}
% TODO: limitations, what each diagnostic licenses you to claim, practical guidance.

% ------------------------------------------------------------------------------------------------------------------------
\section{Conclusion}
% TODO

\bibliographystyle{IEEEtran}
\bibliography{references}

\end{document}
